\chapter{Impact transitoire et modèles de propagateur}\label{ch:mx:transient-impact-and-propagator-models}

Les signes des transactions sont prévisibles pendant des heures~: un achat est suivi d'achats bien après que l'ordre d'un trader donné est terminé. Pourtant, les prix ne suivent pas de tendance.
Quelque chose doit défaire l'impact de chaque transaction, et la vitesse à laquelle cela se produit détermine le coût d'une exécution. Ce chapitre mesure la mémoire du flux d'ordres dans un
marché simulé, ajuste le propagateur qui transforme les transactions en prix, et examine quelles formes de décroissance un modèle d'impact peut avoir sans permettre à un trader
de gagner de l'argent en faisant des allers-retours.

\section{La mémoire longue du flux d'ordres}

Lillo et Farmer (2004) ont montré qu'à la London Stock Exchange l'autocorrélation des signes des ordres décroît approximativement comme une loi de puissance en fonction du décalage, avec un
exposant d'environ 0.6~: le flux d'ordres a une mémoire longue (livre~4, chapitre~17). Lillo, Mike et Farmer (2005) l'ont expliqué par un découpage~: si les métaordres
ont des tailles dont la fonction de répartition cumulée décroît comme $v^{-\alpha}$ et sont exécutés en morceaux égaux, l'autocorrélation des signes décroît comme
$\ell^{-(\alpha-1)}$, soit une mémoire longue dès que $\alpha<2$.

La population \texttt{firm.agentmkt} du chapitre~11 sert de marché au présent chapitre, les signes des preneurs de liquidité étant désormais tirés par séquences de longueur de Pareto d'indice de queue 1.5, comme
si chaque séquence était un métaordre. Sur deux heures simulées, 10\,764 \omterm{def:mx:the-limit-order-book:orders}{ordres au marché} (exécutions regroupées), l'autocorrélation des signes vaut 0.27 au décalage un, 0.106 au
décalage dix et 0.023 au décalage cinquante~; un ajustement en loi de puissance sur les décalages 2 à 50 donne un exposant de 0.62 (\cref{fig:mx:transient-impact-and-propagator-models:kernel},
à gauche), contre 0.5 pour les seules séquences~: les ordres des fondamentalistes, dont les signes suivent le prix, atténuent la mémoire. La réponse
$R(\ell)=\E[\varepsilon_t(p_{t+\ell}-p_t)]$ continue de croître~: 0.52 tick après un ordre, 1.05 après dix, 2.3 après cent.

\section{Le modèle de propagateur}

\begin{definition}[Modèle de propagateur, noyau de décroissance]\label{def:mx:transient-impact-and-propagator-models:propagator}
Le \emph{modèle de propagateur}\index{modèle de propagateur} écrit le prix avant la transaction $t$ comme la somme des impacts de toutes les transactions passées,
\[
p_t=\sum_{s<t}G(t-s)\,\varepsilon_s+\text{bruit},
\]
où le \emph{noyau de décroissance}\index{noyau de décroissance} (le propagateur) $G(\ell)$ est l'impact d'une transaction $\ell$ transactions plus tard, indépendant des autres.
\end{definition}

Bouchaud, Gefen, Potters et Wyart (2004) l'ont proposé et ont montré que le marché se situe en un point critique~: avec des signes corrélés, un $G$ constant ferait
suivre une tendance aux prix et un $G$ à décroissance rapide les ferait revenir à la moyenne~; les prix ne sont diffusifs que lorsque $G$ décroît au rythme qui compense la mémoire des signes.
On estime le noyau à partir des données en écrivant la variation de prix comme $\Delta p_t=\sum_{n\ge0}K(n)\varepsilon_{t-n}$ avec $K(n)=G(n+1)-G(n)$ et en prenant
l'espérance contre les signes passés~: $\E[\varepsilon_{t-k}\Delta p_t]=\sum_nK(n)C(|k-n|)$, un système de Toeplitz résolu par la récursion de Levinson.

\omcode{code/firm/propagator/firm_propagator.py}{40}{48}{Le noyau à partir des corrélations des signes et des corrélations croisées signe--variation de prix, par récursion de Levinson.}\label{lst:mx:transient-impact-and-propagator-models:fit}

Sur le flux simulé, le noyau ajusté vaut 0.541 tick un ordre après une transaction, 0.452 dix ordres plus tard et 0.419 cent plus tard~: il décroît, lentement. Rend-il
les prix diffusifs~? La variance des variations de prix sur $\ell$ ordres, divisée par $\ell$, est constante pour une marche aléatoire
(\cref{fig:mx:transient-impact-and-propagator-models:kernel}, à droite). Pour le marché, elle passe de 1.98 au décalage un à environ 1.35 à partir du décalage dix. Le
propagateur alimenté par les mêmes signes donne 0.29 au décalage un, puis monte à 1.34 à cent~: il explique la variance aux longs décalages mais non le bruit aux courts décalages
(des cotations qui bougent sans transaction), et sa décroissance, tronquée à cent décalages, n'est pas tout à fait assez rapide. Alimenté par les mêmes signes mélangés, sans mémoire, le
même noyau rend les prix sous-diffusifs, de 0.32 à 0.18~: c'est sa décroissance qui compense la mémoire des signes.

\begin{omfigure}
\begin{tikzpicture}
\begin{groupplot}[group style={group size=2 by 1,horizontal sep=1.6cm},width=6.6cm,height=5cm,
  tick label style={font=\footnotesize},label style={font=\small},grid=major,grid style={gray!25},table/col sep=comma,
  xmode=log,xlabel={décalage (ordres)}]
\nextgroupplot[title={\small signes et noyau},ymode=log,ymin=0.01,ymax=1,
  legend columns=1,legend cell align=left,legend style={font=\footnotesize,at={(0.5,-0.3)},anchor=north,draw=none,fill=none}]
\addplot[omDef,very thick,mark=*] table[x=lag,y=acf]{figdata/microstructure/12-transient-impact-and-propagator-models/kernel.csv};
\addplot[omThm,very thick,mark=square*] table[x=lag,y=G]{figdata/microstructure/12-transient-impact-and-propagator-models/kernel.csv};
\legend{autocorrélation des signes $C(\ell)$,noyau $G(\ell)$ (ticks)}
\nextgroupplot[title={\small variance de $p_{t+\ell}-p_t$ selon $\ell$},ymin=0,ymax=2.2,
  legend columns=1,legend cell align=left,legend style={font=\footnotesize,at={(0.5,-0.3)},anchor=north,draw=none,fill=none}]
\addplot[black,very thick,mark=*] table[x=lag,y=market]{figdata/microstructure/12-transient-impact-and-propagator-models/diffusion.csv};
\addplot[omThm,very thick,mark=square*] table[x=lag,y=model]{figdata/microstructure/12-transient-impact-and-propagator-models/diffusion.csv};
\addplot[omProp,very thick,mark=triangle*] table[x=lag,y=shuffled]{figdata/microstructure/12-transient-impact-and-propagator-models/diffusion.csv};
\legend{marché,propagateur sur les signes,propagateur sur les signes mélangés}
\end{groupplot}
\end{tikzpicture}

\omcaption{Le propagateur ajusté sur deux heures simulées. À gauche~: l'autocorrélation des signes (mémoire longue) et le noyau ajusté (décroissance lente). À droite~: la
variance des variations de prix par décalage~; plate signifie diffusif. Données~: \texttt{mx\_prop.kernel\_study}.}\label{fig:mx:transient-impact-and-propagator-models:kernel}
\end{omfigure}

\section{Impact transitoire en temps continu}

\begin{definition}[Modèle d'impact transitoire]\label{def:mx:transient-impact-and-propagator-models:transient}
Un \emph{modèle d'impact transitoire}\index{modèle d'impact transitoire} évalue le coût d'une exécution comme le propagateur en temps continu~: négocier au taux $v(s)$ déplace le prix
à l'instant $t$ de $\int_{s<t}G(t-s)\,v(s)\,ds$, de sorte qu'une exécution $x_1,\dots,x_n$ aux instants $t_1,\dots,t_n$ coûte $\tfrac12x^\top\Gamma x$ avec
$\Gamma_{ij}=G(|t_i-t_j|)$.
\end{definition}

\begin{definition}[Modèle d'Obizhaeva--Wang]\label{def:mx:transient-impact-and-propagator-models:ow}
Le \emph{modèle d'Obizhaeva--Wang}\index{modèle d'Obizhaeva--Wang} est le \omterm{def:mx:transient-impact-and-propagator-models:transient}{modèle d'impact transitoire} d'un carnet d'ordres à forme en bloc et à résilience exponentielle
résilience~: le carnet se reconstitue au taux $\rho$ après avoir été consommé, de sorte que $G(t)=e^{-\rho t}$ à une constante près.
\end{definition}

Obizhaeva et Wang (2013) ont montré que l'exécution optimale dans leur modèle mélange des transactions discrètes et continues~: pour la vente de $X$ sur $[0,T]$, on exécute un bloc
de $X/(\rho T+2)$ au départ, on vend à un taux constant $\rho X/(\rho T+2)$, et on exécute un autre bloc de $X/(\rho T+2)$ à la fin. Le premier bloc utilise le
carnet tel qu'il est~; le négoce continu suit la reconstitution du carnet~; le dernier bloc profite du carnet reconstitué lorsqu'aucune transaction ultérieure ne paiera son impact.
Minimiser $\tfrac12x^\top\Gamma x$ sous $\sum x=X$ donne $x\propto\Gamma^{-1}\mathbf 1$, et sur une grille fine avec un noyau exponentiel, cela reproduit les
deux blocs.

\section{Conditions d'absence de manipulation}

\begin{definition}[Condition d'absence d'arbitrage dynamique]\label{def:mx:transient-impact-and-propagator-models:nda}
Un modèle d'impact vérifie la \emph{condition d'absence d'arbitrage dynamique}\index{condition d'absence d'arbitrage dynamique} lorsqu'aucun aller-retour, une suite d'achats et de ventes
à position nette nulle, n'a un coût espéré négatif~: $x^\top\Gamma x\ge0$ dès que $\sum x=0$ (le noyau est défini positif sur les allers-retours).
\end{definition}

\begin{definition}[Manipulation de prix déclenchée par les transactions]\label{def:mx:transient-impact-and-propagator-models:ttpm}
La \emph{manipulation de prix déclenchée par les transactions}\index{manipulation de prix déclenchée par les transactions} est présente lorsque le moyen le moins coûteux de vendre (ou d'acheter) une
position comprend des transactions dans le sens opposé, bien qu'aucun aller-retour ne soit profitable.
\end{definition}

Huberman et Stanzl (2004) ont montré qu'avec un impact stationnaire dans le temps, seul un impact permanent linéaire exclut le quasi-arbitrage~; Gatheral (2010) a étudié quelles
combinaisons de fonction d'impact et de \omterm{def:mx:transient-impact-and-propagator-models:propagator}{noyau de décroissance} permettent un arbitrage dynamique~; Alfonsi, Schied et Slynko (2012) ont démontré que l'impact doit décroître comme une fonction convexe
non croissante du temps pour exclure à la fois la manipulation standard et la manipulation déclenchée par les transactions. La famille $G(t)=e^{-(t/\tau)^\kappa}$ présente les
trois régimes (\cref{fig:mx:transient-impact-and-propagator-models:family})~: convexe pour $\kappa\le1$~; pour $1<\kappa\le2$, elle est plate, puis décroît, et
reste définie positive (c'est la fonction caractéristique d'une loi stable)~; au-delà de 2, elle décroît trop vite après un début plat et ne l'est pas.

\omcode{code/firm/propagator/firm_propagator.py}{65}{76}{L'aller-retour le plus rentable (la plus petite valeur propre sur le sous-espace des allers-retours) et la liquidation optimale.}\label{lst:mx:transient-impact-and-propagator-models:rt}

Sur 21 transactions réparties sur une unité de temps avec $\tau=0.3$, le meilleur aller-retour perd 0.136 (en unités de l'impact instantané multiplié par le carré de la taille) pour
$\kappa=0.5$, 0.042 pour $\kappa=1$, 0.008 pour 1.5 et 0.002 pour 1.8, et rapporte 0.028 pour $\kappa=2.2$, 0.091 pour 2.5 et 0.213 pour 3. Pour $\kappa=2$, la
matrice du noyau gaussien est singulière sur une grille fine, à la frontière même. Pour $\kappa=1.5$ et 1.8, aucun aller-retour ne rapporte, mais la vente optimale comprend deux
achats, de $-0.086$ et $-0.33$ de l'ordre~: \omterm{def:mx:transient-impact-and-propagator-models:ttpm}{manipulation de prix déclenchée par les transactions}. Avec $\kappa\le1$, il ne vend que.

\begin{omfigure}
\begin{tikzpicture}
\begin{groupplot}[group style={group size=2 by 1,horizontal sep=1.6cm},width=6.6cm,height=5cm,
  tick label style={font=\footnotesize},label style={font=\small},grid=major,grid style={gray!25},table/col sep=comma,
  xlabel={$\kappa$},xmin=0.3,xmax=3.2]
\nextgroupplot[title={\small meilleur aller-retour (profit)},ymin=-0.16,ymax=0.24]
\addplot[omDef,very thick,mark=*] table[x=kappa,y=round_trip]{figdata/microstructure/12-transient-impact-and-propagator-models/family.csv};
\draw[gray,dashed] (axis cs:2,-0.16) -- (axis cs:2,0.24);
\nextgroupplot[title={\small plus petite transaction de la vente optimale},ymin=-0.5,ymax=0.1]
\addplot[omThm,very thick,mark=square*] table[x=kappa,y=min_trade]{figdata/microstructure/12-transient-impact-and-propagator-models/family.csv};
\draw[gray,dashed] (axis cs:1,-0.5) -- (axis cs:1,0.1);
\end{groupplot}
\end{tikzpicture}

\omcaption{Noyaux $e^{-(t/0.3)^\kappa}$ sur 21 transactions. À gauche~: le profit du meilleur aller-retour, positif dès que $\kappa>2$. À droite~: la plus petite transaction de la
vente la moins coûteuse d'une unité, négative (un achat) dès que $\kappa>1$. Données~: \texttt{mx\_prop.kernel\_family}.}\label{fig:mx:transient-impact-and-propagator-models:family}
\end{omfigure}

\section{Estimation du noyau}

L'estimateur de Toeplitz n'a besoin que des signes et des prix, mais ses réponses dépendent de choix. Le retard maximal tronque le noyau, et au-delà, l'ajustement traite
l'impact comme permanent. Regrouper les exécutions en ordres, ou non, change ce qu'est une transaction. D'autres ordres déplacent aussi les prix~: les variations de cotation sans transaction ajoutent du bruit
aux retards courts, comme l'a montré le marché simulé. Et la linéarité du modèle en fonction des signes ignore les tailles~; des extensions laissent le noyau dépendre de la taille de la transaction ou du
fait que la transaction a modifié le prix. Le test de diffusion est le test propre du modèle~: un noyau qui laisse les prix super- ou sous-diffusifs néglige une partie de la
décroissance.

\section{Tutoriel~: le noyau qui payait pour trader}\label{tut:mx:transient-impact-and-propagator-models:tut}

\begin{tutorial}
\textbf{Objectif.} Mesurer la mémoire du flux d'ordres, ajuster un propagateur, vérifier la diffusion et trouver les noyaux qui rémunèrent les allers-retours.
\textbf{État final~:} \cref{fig:mx:transient-impact-and-propagator-models:kernel,fig:mx:transient-impact-and-propagator-models:family}.
\begin{tutsteps}
\item \textbf{Flux.} \texttt{mx\_prop.flow(seed, seconds)}~: signes et milieux des \omterm{def:mx:the-limit-order-book:orders}{ordres au marché} fusionnés provenant de \texttt{firm.agentmkt}, avec des signes à mémoire longue.
\item \textbf{Noyau.} \texttt{firm\_propagator.sign\_acf}, \texttt{response}, \texttt{fit\_kernel}~; \texttt{propagate} pour les prix du modèle.
\item \textbf{Diffusion.} Ratios de variance pour le marché, le modèle et le modèle sur signes mélangés (\texttt{kernel\_study()}).
\item \textbf{Manipulation.} \texttt{impact\_matrix}, \texttt{round\_trip}, \texttt{optimal\_liquidation} sur la famille de noyaux (\texttt{kernel\_family()})~;
\texttt{obizhaeva\_wang(X, rho, T)}~; tracer avec \texttt{fig\_prop.py}.
\end{tutsteps}
\textbf{Pour aller plus loin.} Ajuster le noyau séparément pour les ordres qui ont déplacé le prix et ceux qui ne l'ont pas~; porter la queue de run à 1.8 et observer le
noyau décroître plus vite.
\end{tutorial}

\section{Construction~: le propagateur}\label{bld:mx:transient-impact-and-propagator-models:propagator}

\begin{build}
\bfield{Rôle} Impact transitoire pour les modèles d'exécution des chapitres 14 et 15, les simulations d'impact de marché du chapitre~13, et les
coûts d'inventaire du livre~11.
\bfield{Interface} \texttt{sign\_acf}, \texttt{response}, \texttt{fit\_kernel(eps, p, maxlag)}, \texttt{propagate(eps, G)},
\texttt{impact\_matrix(times, kernel)}, \texttt{round\_trip(Gamma)}, \texttt{optimal\_liquidation(Gamma, X)}, \texttt{obizhaeva\_wang(X, rho, T)}.
\bfield{Règles} Temps en transactions pour le modèle discret~; le système de Toeplitz résolu par \texttt{scipy.linalg.solve\_toeplitz} (Levinson)~; aller-retours sur le
sous-espace $\sum x=0$ de norme unitaire.
\bfield{Tests d'acceptation} \texttt{code/firm/propagator/tests/}~: un noyau en loi de puissance planté, retrouvé à partir de signes persistants~; aller-retours et liquidations
pour un noyau exponentiel, un noyau comprimé et un noyau étiré~; l'optimum discret proche des blocs d'Obizhaeva--Wang.
\bfield{Pour aller plus loin} Noyaux dépendant de la taille et modifiant le prix~; le noyau en temps continu à partir des horodatages~; l'impact transitoire non linéaire de Gatheral.
\end{build}

\begin{omsources}
\item J.-P.~Bouchaud, Y.~Gefen, M.~Potters et M.~Wyart, ``Fluctuations and response in financial markets: the subtle nature of `random' price changes'',
\emph{Quantitative Finance} 4(2), 2004.
\item F.~Lillo et J.~D.~Farmer, ``The long memory of the efficient market'', \emph{Studies in Nonlinear Dynamics and Econometrics} 8(3), 2004.
\item G.~Huberman et W.~Stanzl, ``Price manipulation and quasi-arbitrage'', \emph{Econometrica} 72(4), 2004.
\item F.~Lillo, S.~Mike et J.~D.~Farmer, ``Theory for long memory in supply and demand'', \emph{Physical Review E} 71, 2005.
\item J.~Gatheral, ``No-dynamic-arbitrage and market impact'', \emph{Quantitative Finance} 10(7), 2010.
\item A.~Alfonsi, A.~Schied et A.~Slynko, ``Order book resilience, price manipulation, and the positive portfolio problem'', \emph{SIAM Journal on Financial
Mathematics} 3(1), 2012.
\item A.~A.~Obizhaeva et J.~Wang, ``Optimal trading strategy and supply/demand dynamics'', \emph{Journal of Financial Markets} 16(1), 2013.
\end{omsources}

\section{Exercices}

\begin{exercise}[$\star$]\label{exo:mx:transient-impact-and-propagator-models:1}
L'autocorrélation des signes vaut 0.27 au retard un. Si l'on prédisait que le signe suivant est égal au dernier, à quelle fréquence la prédiction serait-elle juste~?
\end{exercise}

\begin{exercise}[$\star$]\label{exo:mx:transient-impact-and-propagator-models:2}
Pourquoi un noyau constant ferait-il tendre les prix vers une tendance lorsque les signes ont une mémoire longue~?
\end{exercise}

\begin{exercise}[$\star$]\label{exo:mx:transient-impact-and-propagator-models:3}
Dans le \omterm{def:mx:transient-impact-and-propagator-models:ow}{modèle d'Obizhaeva--Wang}, vendre 10\,000 actions sur deux heures avec une résilience $\rho=4$ par heure. Donner les deux blocs et le taux continu.
\end{exercise}

\begin{exercise}[$\star\star$]\label{exo:mx:transient-impact-and-propagator-models:4}
Deux transactions aux instants 0 et 1 avec $G(t)=e^{-t}$. Combien coûte l'aller-retour $(+1,-1)$, et comment vendre une unité avec ces deux transactions~?
\end{exercise}

\begin{exercise}[$\star\star$]\label{exo:mx:transient-impact-and-propagator-models:5}
Expliquer pourquoi le propagateur sur signes mélangés rend les prix sous-diffusifs.
\end{exercise}

\begin{exercise}[$\star\star$]\label{exo:mx:transient-impact-and-propagator-models:6}
Les tailles des métaordres ont une fonction de répartition cumulée décroissante en $v^{-1.5}$ et sont exécutées par morceaux égaux. Quel exposant l'autocorrélation des signes
devrait-elle avoir, et pourquoi la simulation a-t-elle trouvé 0.62~?
\end{exercise}

\begin{exercise}[$\star\star\star$]\label{exo:mx:transient-impact-and-propagator-models:7}
\emph{Programmation.} Pour $\kappa=2.5$, trouver les transactions de l'aller-retour le plus rentable sur la grille de 21 points et décrire sa forme. Pourquoi rapporte-t-il de l'argent~?
\end{exercise}

\begin{exercise}[$\star\star\star$]\label{exo:mx:transient-impact-and-propagator-models:8}
\emph{Trouver la faille.} «~Notre noyau ajusté est défini positif, donc notre modèle d'exécution est exempt de manipulation.~»
\end{exercise}

\section{Problème~: le noyau qui payait pour trader}

\begin{problem}[{Problème du week-end --- le noyau qui payait pour trader}]\label{pb:mx:transient-impact-and-propagator-models:1}
Un quant propose un modèle d'impact dont le noyau reste plat un moment puis chute. Vérifier ce qu'il implique.

\medskip\noindent\textbf{Partie I --- Mémoire.}
\begin{enumerate}
\item Qu'ont trouvé Lillo et Farmer sur le London Stock Exchange~?
\item Énoncer l'explication par découpage et son exposant.
\item Donner l'autocorrélation simulée des signes aux retards 1, 10 et 50, et son exposant ajusté.
\item Donner la réponse après 1, 10 et 100 ordres.
\end{enumerate}

\medskip\noindent\textbf{Partie II --- Le propagateur.}
\begin{enumerate}[resume]
\item Écrire le modèle et le système de Toeplitz associé à son noyau.
\item Donner le noyau ajusté aux retards 1, 10 et 100.
\item Comparer les rapports de variance du marché, du modèle et du modèle sur signes mélangés.
\item Pourquoi le marché est-il critique au sens de Bouchaud et coauteurs~?
\end{enumerate}

\medskip\noindent\textbf{Partie III --- Temps continu.}
\begin{enumerate}[resume]
\item Écrire le coût d'une exécution dans le modèle transitoire.
\item Établir la liquidation optimale $x\propto\Gamma^{-1}\mathbf 1$.
\item Énoncer la solution d'Obizhaeva--Wang et expliquer ses deux blocs.
\item Résoudre l'exercice~3.
\end{enumerate}

\medskip\noindent\textbf{Partie IV --- Manipulation et verdict.}
\begin{enumerate}[resume]
\item Définir l'arbitrage dynamique et la manipulation déclenchée par les transactions.
\item Énoncer la condition d'Alfonsi, Schied et Slynko.
\item Donner le meilleur aller-retour pour $\kappa=0.5$, 1, 1.5, 1.8, 2.2, 2.5 et 3.
\item Quels noyaux font-ils comporter des achats à la vente optimale~?
\item Que se passe-t-il exactement pour $\kappa=2$~?
\item Énoncer le \emph{résultat nommé}~: le profit d'aller-retour sous un noyau qui décroît trop vite après un départ plat, et le paramètre de décroissance pour lequel il s'annule.
\item Dans quel régime se situe le noyau proposé~?
\item En une phrase~: à quoi doit ressembler un noyau pour être sûr~?
\end{enumerate}
\end{problem}

\section{Questions d'entretien}

\begin{interviewq}[$\star$ \iqroles{chercheur}]\label{iq:mx:transient-impact-and-propagator-models:1}
Les signes des ordres sont prévisibles, mais pas les prix. Comment est-ce possible~?
\end{interviewq}

\begin{interviewq}[$\star\star$ \iqroles{chercheur}]\label{iq:mx:transient-impact-and-propagator-models:2}
Comment estimeriez-vous le noyau d'impact à partir des transactions et des cotations~?
\end{interviewq}

\begin{interviewq}[$\star\star$ \iqroles{chercheur}]\label{iq:mx:transient-impact-and-propagator-models:3}
Quelle condition un modèle d'impact doit-il satisfaire pour exclure les allers-retours profitables~?
\end{interviewq}

\begin{interviewq}[$\star\star$ \iqroles{trader}]\label{iq:mx:transient-impact-and-propagator-models:4}
Pourquoi l'exécution optimale dans le \omterm{def:mx:transient-impact-and-propagator-models:ow}{modèle d'Obizhaeva--Wang} commence-t-elle et se termine-t-elle par des blocs~?
\end{interviewq}

\begin{interviewq}[$\star\star$ \iqroles{développeur}]\label{iq:mx:transient-impact-and-propagator-models:5}
Calculez le coût d'un calendrier de $n$ transactions sous un noyau transitoire. Quelle est la complexité, et comment l'accéléreriez-vous~?
\end{interviewq}

\begin{interviewq}[$\star\star\star$ \iqroles{chercheur}]\label{iq:mx:transient-impact-and-propagator-models:6}
Votre optimiseur, alimenté par un noyau ajusté, indique qu'une liquidation doit acheter pendant un temps. Est-ce un bogue~?
\end{interviewq}
